\documentclass{beamer}
\usetheme{Madrid}

% Standard packages
\usepackage[utf8]{inputenc}
\usepackage{booktabs}
\usepackage{graphicx}

% Presentation Info
\title[Low-Latency DASH Streaming]{Low-Latency DASH Live Streaming Pipeline}
\subtitle{CS 418 Project 3 - Final Presentation}
\author[Atakan Karata\c{s}]{Atakan Karata\c{s}}
\institute[Özye\v{g}in University]{
    Department of Computer Engineering\\
    \"{O}zye\v{g}in University, Istanbul, Turkey
}
\date[May 2026]{May 30, 2026}

\begin{document}

% Slide 1: Title Page
\begin{frame}
    \titlepage
\end{frame}

% Slide 2: Outline
\begin{frame}{Outline}
    \tableofcontents
\end{frame}

\section{Project Goal \& Motivation}

% Slide 3: Goal & Motivation
\begin{frame}{Project Goal \& Motivation}
    \begin{block}{Core Objective}
        Build a low-latency live streaming pipeline to achieve glass-to-glass latency of \textbf{under 5 seconds} using open-source tools.
    \end{block}
    \vspace{0.2cm}
    \begin{itemize}
        \item \textbf{Use Cases:} Real-time interaction, remote monitoring, interactive live shows.
        \item \textbf{Challenges:} Standard HTTP streaming protocols (DASH/HLS) introduce 6-30 seconds of delay due to segment buffering.
        \item \textbf{Key Approach:} Low-latency DASH (LL-DASH) using chunked transfer encoding.
    \end{itemize}
\end{frame}

\section{System Architecture}

% Slide 4: System Architecture
\begin{frame}{End-to-End Pipeline Architecture}
    The streaming delivery chain consists of five main components:
    \vspace{0.2cm}
    \begin{itemize}
        \item \textbf{OBS Studio:} Media source. Captures video and burns high-precision wall-clock time overlay.
        \item \textbf{RTMP Link:} Encoded media transmission from OBS to local packaging process.
        \item \textbf{Modified FFmpeg:} Packages incoming RTMP stream into low-latency DASH chunks/segments.
        \item \textbf{NGINX \& node-gpac-dash:} Serves the dash.js player and delivers segments progressively.
        \item \textbf{dash.js Player:} Browser playback engine with live catchup.
    \end{itemize}
\end{frame}

\section{FFmpeg Source Build}

% Slide 5: FFmpeg Source Modification (Appendix A)
\begin{frame}{FFmpeg Source Modification}
    \begin{block}{The Problem}
        FFmpeg normally writes DASH segments with a temporary \texttt{.tmp} extension, renaming them upon completion. This blocks node-gpac-dash from streaming files mid-write.
    \end{block}
    \vspace{0.2cm}
    \begin{block}{The Solution (Appendix A)}
        Patch \texttt{libavformat/dashenc.c} to write directly to the final file name:
        \vspace{0.1cm}
        \small
        \begin{description}
            \item[Original:] \texttt{use\_rename ? "\%s.tmp" : "\%s"}
            \item[Modified:] \texttt{use\_rename ? "\%s" : "\%s"}
        \end{description}
    \end{block}
    \vspace{0.1cm}
    \begin{itemize}
        \item Compiled with flags: \texttt{--enable-demuxer=dash --enable-libxml2}.
    \end{itemize}
\end{frame}

\section{Core Validation}

% Slide 6: Core Validation
\begin{frame}{Core Validation}
    \begin{itemize}
        \item \textbf{HTTP Chunked Encoding:} Verified via Network inspector showing \texttt{Transfer-Encoding: chunked} for DASH segment responses.
        \item \textbf{Stabilizing Encoder Lag:} High OBS encoding resolutions caused CPU lag. Resolved by reducing settings to:
        \begin{itemize}
            \item Resolution: \textbf{852x480}
            \item Bitrate: \textbf{1200 kbps} (30 FPS)
        \end{itemize}
        \item \textbf{Baseline Result:} 1-second segment profile reached stable \textbf{~4.0 seconds} of glass-to-glass latency without stalls.
    \end{itemize}
\end{frame}

\section{Experiments}

% Slide 7: Segment Duration Experiment
\begin{frame}{Segment Duration Experiment}
    Varying segment duration (\texttt{seg\_duration}) and keyframe interval (\texttt{keyint}) at 30 FPS:
    \vspace{0.2cm}
    \begin{table}
        \centering
        \scriptsize
        \begin{tabular}{cccr}
            \toprule
            \textbf{Seg. Duration (s)} & \textbf{Keyint} & \textbf{Latency (ms)} & \textbf{Stability} \\
            \midrule
            \textbf{1} & \textbf{30} & \textbf{4000} & \textbf{Stable (Target Met)} \\
            2 & 60 & 5500 & Stable \\
            3 & 90 & 5500 & Stable \\
            4 & 120 & 5500 & Stable \\
            5 & 150 & 7000 & Stalls \\
            6 & 180 & 7500 & Stable but slow \\
            \bottomrule
        \end{tabular}
    \end{table}
    \begin{itemize}
        \item \textbf{Insight:} Latency scales with segment size. Only the 1s segment profile achieved the required under-5s latency.
    \end{itemize}
\end{frame}

% Slide 8: Fragment Duration Experiment
\begin{frame}{Fragment Duration Experiment}
    Testing fragment duration (\texttt{frag\_duration}) with fixed 4s segments and \texttt{keyint=120}:
    \vspace{0.1cm}
    \begin{table}
        \centering
        \tiny
        \begin{tabular}{ccr}
            \toprule
            \textbf{Fragment Duration (s)} & \textbf{Latency (ms)} & \textbf{Stability} \\
            \midrule
            0.033 & 5000 & Stable \\
            0.066 & 5500 & Stable \\
            0.100 & 5500 & Stable \\
            0.200 & 5500 & Stable \\
            0.500 & 6000 & Stable \\
            1.000 & 6000 & Minor waiting / catchup \\
            2.000 & 6000 & Gap Controller adjustments \\
            4.000 & 9000 & Instability / stalls \\
            \bottomrule
        \end{tabular}
    \end{table}
    \vspace{-0.1cm}
    \begin{itemize}
        \item \textbf{Insight:} Fragment size alone does not bypass segment-level buffer delays. Large fragments cause repeated player stalls.
    \end{itemize}
\end{frame}

\section{Bonus: QR Latency Monitor}

% Slide 9: Bonus: QR Latency Monitor
\begin{frame}{Bonus: QR Latency Monitor}
    \begin{itemize}
        \item \textbf{System Clock Overlay:} Generates QR code containing local clock timestamp.
        \item \textbf{Payload format:} \texttt{cs418-project3-ts=<ISO timestamp>}.
        \item \textbf{Player-side Decoding:} Samples frames from HTML5 video element into an offscreen canvas.
        \item \textbf{Real-Time Latency:} Uses the browser-based \texttt{jsQR} library to decode the QR code and subtract it from the player clock.
        \item \textbf{Advantage:} Enables automated, sub-second precision latency monitoring without manual stopwatch screenshots.
    \end{itemize}
\end{frame}

\section{Conclusion}

% Slide 10: Conclusion
\begin{frame}{Conclusion}
    \begin{itemize}
        \item \textbf{Goal Achieved:} Under 5 seconds latency is achieved at \textbf{4.0 seconds} adjusted latency.
        \item \textbf{System Sensitivity:} Low latency requires tuning the entire delivery chain (encoder settings, proxy buffering, segment/chunk sizes).
        \item \textbf{Chunked Delivery:} Progressive HTTP/1.1 chunked transfer is mandatory to prevent segment buffering delays.
        \item \textbf{Verification:} All test suites, formatting, and validation scripts pass locally.
    \end{itemize}
\end{frame}

\section*{Thank You}

% Slide 11: Thank You
\begin{frame}{Thank You}
    \centering
    \vspace{1.0cm}
    \Huge Questions?
    \vspace{0.5cm}
    \Large \\ Q&A Session
\end{frame}

\end{document}

